\section{Analysis}
\label{sec:mechanism}

Section~\ref{sec:results} is an existence proof over $17{,}712$ rules, two seen
and two unseen models, and both splits. This section explains \emph{why} it holds
and why it generalizes beyond the exact rules we searched: the accuracy cost of
stopping on consensus is a directional property of the reasoning trajectory
itself, not of a threshold or of probe density.

\subsection{Stopping destroys more than it saves: a $35{:}1$ effect}
Consider only the stops where committing early \emph{changes} correctness relative
to the trajectory's own final answer, and count two outcomes: (final-correct,
stop-wrong)---a correct recovery destroyed---versus (final-wrong,
stop-correct)---a wrong answer luckily banked. Pure sampling noise predicts a
$\approx1{:}1$ ratio; a genuine ``continued reasoning corrects'' effect predicts
$\gg1$. On the frozen replay this ratio is $\mathbf{35{:}1}$ for the naive
consensus stopper ($1055$ destroyed vs.\ $30$ banked; $34.3{:}1$ and $35.9{:}1$ on
\dsseven{} and \qweneight{} separately) and $15$--$18{:}1$ for the more
conservative variants---all far above $1{:}1$. Stopping on consensus removes a
correct-in-the-end answer about $35$ times more often than it rescues a wrong one.
This is the recovery phenomenon of \S\ref{sec:exp-fc} ($65.5\%$, $76.3\%$), now at
the main $16$K/$32$K budgets and as a signed effect: the loss is real and
directional, not sampling variance.

\subsection{The accuracy tax is intrinsic to the trajectory}
Why is the ratio structural? For a rule stopping at token position $s$, the
accuracy drop depends only on $s$: committing at $s$ forgoes whatever correction
the trajectory would have made after it. It is a property of the reasoning, not of
the probe, so it survives \emph{any} probing scheme. This is what never reaches
zero in the sweep. The probe grid (intervals $64/128/256$ tokens) makes
essentially every stop position reachable and the observed floor is the minimum
over these schedules; a sparser schedule can only stop later or at a subset of
positions, never at a safer one the grid missed. Hence no threshold on the
consensus signal---in or beyond our schema---escapes the accuracy tax; the
searched space is a lower bound on the cost, not a lucky miss.

\subsection{Ruling out the ``it is just dense probing'' objection}
The natural rebuttal is that the negative result is an artifact of charging dense
re-probing. It is not, and separating two costs shows why. Besides the accuracy
tax, a rule pays a \emph{probe tax}: the probe output $p$ (Section~\ref{sec:accounting}).
Under dense re-probing a late-stopping rule pays $p$ over a long prefix, which
turns the \emph{net} savings of the safe (entropy) family negative even though its
\emph{gross} savings is positive (Table~\ref{tab:grossnet}). Unlike the accuracy
tax, the probe tax scales with probe density and is reducible (sparser or shorter
probes, KV-cache reuse). The two are separable: the \emph{savings} collapse is
partly a probe-tax artifact and is addressable, but the \emph{accuracy} floor is
probe-independent and is what makes the corner empty.

\begin{table}[t]
  \centering\small
  \begin{tabular}{lcc}
    \toprule
    Operating point & Gross saving & Net saving \\
                    & $(B-s)/B$    & $(B-s-p)/B$ \\
    \midrule
    Ours (target conservative)  & $+0.6\%$  & $-4.0\%$ \\
    Ours (target balanced)      & $+10.0\%$ & $+7.8\%$ \\
    Ours (target token\_efficient) & $+21.5\%$ & $+19.6\%$ \\
    \bottomrule
  \end{tabular}
  \caption{Gross vs.\ net savings; the gap is the probe tax. The
  target-conservative point's negative net saving is a probing-density artifact,
  not evidence that stopping wastes tokens---whereas the accuracy tax above it is
  intrinsic.}
  \label{tab:grossnet}
\end{table}

\subsection{It is the signal, not early exit}
The reproductions of \S\ref{sec:exp-prior} localize the cause. Every method that
stops on \emph{probe agreement}---the swept schema and CertaIndex alike---either
fails the safety bar or collapses outright; CertaIndex's $70\pp$ loss at $99.8\%$
stop rate is the accuracy tax paid in full. The lone method that stays near full
accuracy, DEER, does not read probe agreement at all: it stops on the model's own
confidence in a trial answer at a reasoning boundary, and in
Figure~\ref{fig:pareto} it sits \emph{above} the consensus Pareto boundary at the
same accuracy drop. This is direct evidence that the failure is a property of the
\emph{consensus signal}, not of early exit. (DEER's edge is model-dependent,
$+0.8\pp$ on \qweneight{} vs.\ $-4.8\pp$ on \dsseven{}, and these are
frozen-trajectory reproductions.)

\subsection{Implication: change the signal, not the threshold}
Much prior effort tunes consensus-based stops
\citep{dynasor_certaindex,aggarwal2023adaptive,li2024esc}; our frontier shows the
ceiling of that approach is low, and the mechanism above says no amount of tuning
raises it. The constructive implication is that a safe signal should be about
whether the \emph{trajectory itself} will still change, not about the agreement of
backward-looking probes. As a first step, we sketch in Appendix~\ref{sec:boundary}
an online controller built on DEER's boundary-confidence signal (fast-commit on
high confidence, else a verification branch) that recovers substantial savings at
near-neutral accuracy on dev; we stress it is exploratory (three seeds,
seed-sensitive, outside the preregistration) and report it only to show that
stepping outside the consensus-signal space escapes the accuracy tax we proved
unavoidable inside it. Verifier- and process-reward-based termination
\citep{lightman2024verify,wang2024mathshepherd} and learned halting
\citep{graves2016act,schuster2022confident} point the same way. Turning any of
these into a \emph{confirmed} method---with its own train/dev$\rightarrow$test
validation---is the natural next paper; our contribution is to establish
rigorously which signal not to build it on. Two caveats bound the claim: it
concerns \emph{consensus-based} rules within the searched schema and dense-probe
regime, and it does not say early exit is impossible---only that this signal
cannot make it safe.
